
\documentclass{article}
\usepackage{colortbl}
\usepackage{makecell}
\usepackage{multirow}
\usepackage{supertabular}

\begin{document}

\newcounter{utterance}

\centering \large Interaction Transcript for game `eqbench', experiment `eq\_bench', episode 86 with reflexion{-}llama{-}3.1{-}8b{-}instruct.
\vspace{24pt}

{ \footnotesize  \setcounter{utterance}{1}
\setlength{\tabcolsep}{0pt}
\begin{supertabular}{c@{$\;$}|p{.15\linewidth}@{}p{.15\linewidth}p{.15\linewidth}p{.15\linewidth}p{.15\linewidth}p{.15\linewidth}}
    \# & \multicolumn{2}{c}{Player} && \multicolumn{2}{c}{Game Master} \\
    \hline

    \theutterance \stepcounter{utterance}  
    & & & \multicolumn{4}{p{0.6\linewidth}}{
        \cellcolor[rgb]{0.9,0.9,0.9}{
            \makecell[{{p{\linewidth}}}]{
                \texttt{\tiny{[P1$\langle$GM]}}
                \texttt{You have attempted this same task before and did not succeed. These are the plans you wrote after those attempts.} \\
\texttt{Trial \#1: New plan: I was too focused on the surface{-}level banter between Imogen and Sebastian, but I didn't consider the underlying power dynamics and Imogen's potential emotional state. Next time, I will pay closer attention to the tone and language used by Sebastian, which seems to be condescending and dismissive. I will also consider Imogen's initial reaction to seeing Sebastian, which suggests she may be feeling some level of unease or discomfort. I will score the emotions based on the dialogue and Sebastian's behavior, with a focus on Imogen's likely emotional response to his condescending tone. } \\
\\ 
\texttt{Confused: 0} \\
\texttt{Annoyed: 9} \\
\texttt{Resigned: 6} \\
\texttt{Guilty: 0} \\
\\ 
\\ 
\\ 
\texttt{Your task is to predict the likely emotional responses of a character in this dialogue:} \\
\\ 
\texttt{Imogen: Sebastian, I wasn't expecting to see you here.} \\
\texttt{Sebastian: The feeling's mutual, Imogen.} \\
\texttt{Imogen: It's been a while.} \\
\texttt{Sebastian: Not nearly long enough, apparently.} \\
\texttt{Imogen: I see you're still as charming as ever.} \\
\texttt{Sebastian: And you're still as naive.} \\
\texttt{Imogen: That's one way to put it. I prefer to think of it as principled.} \\
\texttt{Sebastian: Principles don't pay the bills, Imogen.} \\
\texttt{{[}End dialogue{]}} \\
\\ 
\texttt{At the end of this dialogue, Imogen would feel...} \\
\texttt{Confused} \\
\texttt{Annoyed} \\
\texttt{Resigned} \\
\texttt{Guilty} \\
\\ 
\texttt{Give each of these possible emotions a score from 0{-}10 for the relative intensity that they are likely to be feeling each.} \\
\\ 
\texttt{You must output in the following format, including headings (of course, you should give your own scores), with no additional commentary:} \\
\\ 
\texttt{Confused: <score>} \\
\texttt{Annoyed: <score>} \\
\texttt{Resigned: <score>} \\
\texttt{Guilty: <score>} \\
\\ 
\\ 
\texttt{{[}End of answer{]}} \\
\\ 
\texttt{Remember: zero is a valid score, meaning they are likely not feeling that emotion. You must score at least one emotion > 0.} \\
\\ 
\texttt{Your answer:} \\
            }
        }
    }
     \\ \\

    \theutterance \stepcounter{utterance}  
    & \multicolumn{4}{p{0.6\linewidth}}{
        \cellcolor[rgb]{0.561,0.737,0.561}{
            \makecell[{{p{\linewidth}}}]{
                \texttt{\tiny{[P1$\rangle$GM]}}
                \texttt{Confused: 0} \\
\texttt{Annoyed: 9} \\
\texttt{Resigned: 2} \\
\texttt{Guilty: 0} \\
            }
        }
    }
    & & \\ \\

    \theutterance \stepcounter{utterance}  
    & & & \multicolumn{2}{p{0.3\linewidth}}{
        \cellcolor[rgb]{0.9,0.9,0.9}{
            \makecell[{{p{\linewidth}}}]{
                \texttt{\tiny{[GM$|$GM]}}
                \texttt{Confused: 0} \\
\texttt{Annoyed: 9} \\
\texttt{Resigned: 2} \\
\texttt{Guilty: 0} \\
            }
        }
    }
    & & \\ \\

    \theutterance \stepcounter{utterance}  
    & & & \multicolumn{2}{p{0.3\linewidth}}{
        \cellcolor[rgb]{0.9,0.9,0.9}{
            \makecell[{{p{\linewidth}}}]{
                \texttt{\tiny{[GM$|$GM]}}
                \texttt{Confused: 0} \\
\texttt{Annoyed: 8} \\
\texttt{Resigned: 0} \\
\texttt{Guilty: 0} \\
            }
        }
    }
    & & \\ \\

    \theutterance \stepcounter{utterance}  
    & & & \multicolumn{2}{p{0.3\linewidth}}{
        \cellcolor[rgb]{0.9,0.9,0.9}{
            \makecell[{{p{\linewidth}}}]{
                \texttt{\tiny{[GM$|$GM]}}
                \texttt{game\_result = LOSE} \\
            }
        }
    }
    & & \\ \\

    \theutterance \stepcounter{utterance}  
    & & \multicolumn{4}{p{0.6\linewidth}}{
        \cellcolor[rgb]{0.95,0.95,0.95}{
            \makecell[{{p{\linewidth}}}]{
                \texttt{\tiny{[P1$\leftrightarrow$P]}}
                \texttt{Post{-}episode Reflexion (not sent to GM):} \\
\texttt{New plan: In my previous attempts, I focused on the surface{-}level emotions expressed by Imogen and Sebastian, but I didn't consider the underlying power dynamics and Imogen's potential emotional state. I also didn't pay enough attention to the tone and language used by Sebastian, which seems to be condescending and dismissive. Next time, I will take a more nuanced approach by analyzing the dialogue from Imogen's perspective, considering her initial reaction to seeing Sebastian and her responses to his condescending tone. I will also examine the context of the conversation, including the fact that Sebastian is making a comment about Imogen's financial situation, which may be a sensitive topic for her. Based on this analysis, I will reassess the emotions and adjust the scores accordingly. } \\
\\ 
\texttt{Confused: 0} \\
\texttt{Annoyed: 10} \\
\texttt{Resigned: 8} \\
\texttt{Guilty: 0} \\
            }
        }
    }
    & \\ \\

\end{supertabular}
}

\end{document}
